\section*{第 \ref{ch:b2:linalg} 章 --- 线性代数}

\begin{solution}{exo:b2:linalg:1}
三个线性型位于 $3$ 维对偶空间中：只需证明它们线性无关。若关系式
$\alpha\varphi_1 + \beta\varphi_2 + \gamma\varphi_3 = 0$ 在
$(1,0,0), (0,1,0), (0,0,1)$ 处求值，得 $\alpha + \gamma = 0$，$\alpha +
\beta = 0$，$\beta + \gamma = 0$，由此 $\alpha = \beta = \gamma =
0$。

求基 $(u_1, u_2, u_3)$：解方程 $\varphi_i(u_j) =
\delta_{ij}$。设 $u_j = (x, y, z)$，对 $u_1$ 有 $x + y = 1$，
$y + z = 0$，$x + z = 0$，得 $u_1 = \bigl(\tfrac12, \tfrac12,
-\tfrac12\bigr)$；同理 $u_2 = \bigl(-\tfrac12, \tfrac12,
\tfrac12\bigr)$，$u_3 = \bigl(\tfrac12, -\tfrac12, \tfrac12\bigr)$。
\end{solution}

\begin{solution}{exo:b2:linalg:2}
第一个矩阵：乘积非零的唯一置换把 $1 \mapsto
3$，$2 \mapsto 2$，$3 \mapsto 1$ --- 即对换 $(1\,3)$，
符号为 $-1$，行列式为 $-abc$。

第二个矩阵：乘积非零的置换不能混合两个分块，
（连接二者的元素为 $0$），因此它分解为
$\{1,2\}$ 上的一个置换与 $\{3,4\}$ 上的一个置换之积，其符号为
两者符号之积，于是和式分解为
\[
(ad - bc)(eh - fg) .
\]
由此可得一般规律（证明相同，结论同样成立）：分块对角矩阵的
行列式等于各分块行列式的
乘积。
\end{solution}

\begin{solution}{exo:b2:linalg:3}
$F$ 由两个独立方程 $\varphi_1(x,y,z,t)
= x + y - z - t = 0$ 与 $\varphi_2 = x - 2y = 0$ 确定。由
\cref{thm:b2:linalg:annihilator} 反用可得 $F^\circ =
\operatorname{Vect}(\varphi_1, \varphi_2)$ --- 二者由构造属于 $F^\circ$，
它们线性无关（互不成比例），且 $\dim F^\circ
= 4 - \dim F = 4 - 2 = 2$，因 $\dim F = 2$（两个独立
方程位于 $\R^4$ 中）。基为 $(\varphi_1, \varphi_2)$；维数：
$2 + 2 = 4$，符合定理要求。
\end{solution}

\begin{solution}{exo:b2:linalg:4}
$\varphi_i$ 是 $n + 1$ 个线性型，位于 $(n+1)$ 维空间上：
只需证明线性无关。若 $\sum_i \lambda_i \varphi_i = 0$，将其
在节点的拉格朗日多项式 $L_j$ 上求值，得 $\lambda_j = 0$。
所求基为 $(L_0, \dots, L_n)$，因为 $\varphi_i(L_j) =
L_j(a_i) = \delta_{ij}$。

对于节点为 $0, \frac12, 1$ 的积分线性型，定义在 $\R_2[X]$ 上：
$\int_0^1 P = \sum_i c_i P(a_i)$，其中 $c_i = \int_0^1 L_i$。计算：
$L_0 = 2(X - \tfrac12)(X - 1)$，$\int_0^1 L_0 = \frac16$；$L_1 =
-4X(X-1)$，$\int_0^1 L_1 = \frac46$；$L_2 = 2X(X - \tfrac12)$，
$\int_0^1 L_2 = \frac16$。由此得
\[
\int_0^1 P = \frac{1}{6}\Bigl(P(0) + 4P\bigl(\tfrac12\bigr) +
P(1)\Bigr)
\quad (P \in \R_2[X]) :
\]
即辛普森公式，对二次多项式精确成立 --- 这是
关于对偶基的一则结论。
\end{solution}

\begin{solution}{exo:b2:linalg:5}
$\operatorname{im} u = Ka$，取 $a \neq 0$；则 $u(x) =
\varphi(x)\,a$，其中 $\varphi(x)$ 为 $u(x)$ 在 $a$ 上的坐标
--- 关于 $x$ 线性。求迹：将 $a = e_1$ 扩充为一组基，
$u$ 的矩阵各列为 $\varphi(e_j)\,e_1$，故其唯一的对角元
为 $\varphi(e_1) = \varphi(a)$，即 $\operatorname{tr} u =
\varphi(a)$。于是
\[
u^2(x) = \varphi(x)\, u(a) = \varphi(x)\varphi(a)\, a
= (\operatorname{tr} u)\, u(x).
\]
行列式分两种情形计算。\emph{若 $\varphi(a) \neq 0$：}任取
超平面 $\ker\varphi$ 的一组基，并添上 $a$。于是 $u$ 把
$\ker\varphi$ 映为零（此时 $u(x) = \varphi(x)a = 0$），且 $u(a) =
\varphi(a)\,a$，故 $I + u$ 的矩阵为对角阵 $(1, \dots, 1,\,
1 + \varphi(a))$，因此 $\det(I + u) = 1 + \varphi(a) = 1 +
\operatorname{tr} u$。\emph{若 $\varphi(a) = 0$：}则 $a \in
\ker\varphi$；取 $\ker\varphi$ 的一组基，其第一个向量为
$a$，并添上向量 $b$，满足 $\varphi(b) = 1$。于是 $I + u$
固定 $\ker\varphi$ 的基，并把 $b \mapsto b + a$：
为三角阵，对角元均为 1，$\det(I + u) = 1 = 1 +
\operatorname{tr} u$。两种情形均与该公式相符。
\end{solution}

\begin{solution}{exo:b2:linalg:6}
由\cref{exo:b2:linalg:9}可知，超平面为 $H_A = \{M :
\operatorname{tr}(AM) = 0\}$，其中 $A \neq 0$。

\emph{若 $A = \lambda I$：}$H_A$ 为迹为零的超平面；
$n$-轮换置换的矩阵（在位置 $(i, i+1)$
及 $(n, 1)$ 处为 1）可逆（其行列式为 $\pm 1$，由
\cref{ex:b2:linalg:permexample}的计算），且迹为零。

\emph{若 $A$ 不是数量矩阵：}先求可逆矩阵 $P$，使得
$B = P^{-1}AP$ 有一个非零的非对角元 $b_{ji}$（$j \neq
i$）。事实上，若 $A$ 本身已有此类元素，取 $P = I$；若 $A$
为对角阵且有两个不同的对角元 $d_1 \neq d_2$，则
以倍加矩阵 $P = I + E_{12}$ 作共轭，得非对角元
$d_1 - d_2 \neq 0$（计算：$P^{-1}AP = A + (d_1 - d_2)E_{12}$）；
而对角元全相等的对角阵为
数量矩阵，已排除。现取 $M' = I + tE_{ij}$，其中 $t =
-\operatorname{tr}(B)/b_{ji}$，则
\[
\operatorname{tr}(BM') = \operatorname{tr} B + t\,b_{ji} = 0,
\]
且 $M'$ 可逆（三角阵，对角元均为 1）。还原
共轭后，$M = PM'P^{-1}$ 可逆，且
$\operatorname{tr}(AM) = \operatorname{tr}(BM') = 0$，故 $M \in H_A$。
\end{solution}

\begin{solution}{exo:b2:linalg:7}
由置换公式，$\det(I + tA)$ 是关于 $t$ 的多项式；
其常数项为 $1$（$t = 0$）。求 $t$ 的系数时，将
$\det$ 按列 $e_j + t\,c_j(A)$ 作交错型展开；由
多重线性，$t$ 的一次项恰把某一个 $e_j$ 换成
$c_j(A)$：
\[
\sum_{j} \det(e_1, \dots, c_j(A), \dots, e_n)
= \sum_j a_{jj} = \operatorname{tr} A ,
\]
（其余各列均为标准基向量，$c_j(A)$ 位于
第 $j$ 个位置的行列式取出第 $j$ 个对角元）。因此导数在 $0$ 处为
$\operatorname{tr} A$。

设 $g(t) = \det(\eu^{tA})$。由群性质，$g(s + t) =
g(s)g(t)$（$\det$ 的乘性），$g$ 可微，且
$g'(0) = \operatorname{tr} A$（由上式，$\eu^{tA} = I + tA +
O(t^2)$）。可微同态 $(\R, +) \to (\R^*, \times)$
满足 $g' = g'(0)\,g$（对 $g(s+t)$ 关于 $s$ 在 $0$ 处求导），故
$g(t) = \eu^{t\operatorname{tr} A}$，由
$y' = cy$ 满足 $y(0) = 1$ 的解的唯一性（本科一年级分册）。
\end{solution}

\begin{solution}{exo:b2:linalg:8}
设 $v_k = (1, j^k, j^{2k})^{\mathsf T}$，$k = 0, 1, 2$。利用 $1
+ j + j^2 = 0$ 与 $j^3 = 1$：
\[
C v_k =
\begin{pmatrix}
a + b j^k + c j^{2k}\\
c + a j^k + b j^{2k}\\
b + c j^k + a j^{2k}
\end{pmatrix}
= (a + b j^k + c j^{2k})
\begin{pmatrix} 1\\ j^k\\ j^{2k}\end{pmatrix},
\]
（验证第二行：$j^k(a + bj^k + cj^{2k}) = aj^k + bj^{2k} +
cj^{3k} = c + aj^k + bj^{2k}$）。故 $v_k$ 为特征向量，
特征值为 $\lambda_k = a + bj^k + cj^{2k}$。$v_k$ 构成一组基
（$1, j, j^2$ 互异，其范德蒙矩阵可逆），故 $C$ 可对角化，
并有
\[
\det C = \lambda_0\lambda_1\lambda_2
= (a+b+c)(a + bj + cj^2)(a + bj^2 + cj).
\]
\end{solution}

\begin{solution}{exo:b2:linalg:9}
映射 $\Theta \colon A \mapsto \operatorname{tr}(A\,\cdot)$ 是
从 $\mathcal{M}_n(K)$ 到其对偶空间的线性映射，二者维数相等，
均为 $n^2$，只需证明单射。若 $\operatorname{tr}(AM) =
0$ 对所有 $M$ 成立，取 $M = E_{ji}$，则 $\operatorname{tr}(A E_{ji}) =
a_{ij} = 0$ 对所有 $i, j$ 成立，于是 $A = 0$，故 $\Theta$ 为同构。

迹的唯一性（\cref{prop:b2:linalg:trace}）：设线性型 $t$
使所有换位子为零，则它是 $\operatorname{tr}(A\,\cdot)$，
对某个 $A$ 满足 $\operatorname{tr}(A(MN - NM)) = 0$ 对所有 $M, N$ 成立，即
$\operatorname{tr}((AM - MA)N) = 0$ 对所有 $N$ 成立（由迹的轮换性），即
$AM = MA$ 对所有 $M$ 成立（由 $\Theta$ 的单射性）：$A$ 与一切矩阵
可交换，因而为数量矩阵（$A$ 与所有 $E_{ij}$ 可交换，
迫使非对角元为 $0$、对角元彼此相等），故 $t = c
\operatorname{tr}$。
\end{solution}

\begin{solution}{exo:b2:linalg:10}
在 $\C$ 上讨论（实矩阵幂零当且仅当其作为复矩阵幂零，
幂零性即 $u^n = 0$）。

\emph{第一步：$\operatorname{tr}(u^k) = 0$（$k \geq 1$）。}由
归纳法，$u^k v - v u^k = k\, u^k$：当 $k = 1$ 时即为
假设；归纳步为
\[
u^{k+1}v - vu^{k+1} = u^k(uv - vu) + (u^k v - v u^k)u
= u^{k+1} + k\,u^{k+1} .
\]
取迹得：$0 = \operatorname{tr}(u^k v) -
\operatorname{tr}(vu^k) = k \operatorname{tr}(u^k)$，故
$\operatorname{tr}(u^k) = 0$。

\emph{第二步：（在 $\C$ 上）幂迹全为零蕴含幂零。}设
$\lambda_1, \dots, \lambda_r$ 为互异的非零特征值，
$u$ 的重数分别为 $m_1, \dots, m_r$（在特征多项式中计算，
该多项式在 $\C$ 上完全分裂 ---
\cref{ch:b2:reduction}）。幂迹为 $\operatorname{tr}(u^k) =
\sum_i m_i \lambda_i^k$（三角化后，三角阵的
$k$ 次幂的对角元恰为对角元的 $k$ 次幂）。方程组 $\sum_i m_i
\lambda_i^k = 0$（$k = 1, \dots, r$）关于未知量 $m_i\lambda_i$ 是可逆
的范德蒙型方程组（系数矩阵为 $(\lambda_i^{k-1})$ 与
对角阵 $\lambda_i$ 之积，所有 $\lambda_i \neq 0$ 且互异）：每个 $m_i
\lambda_i = 0$，这在 $m_i \geq 1$ 时不可能，除非 $r = 0$。于是 $u$
没有非零特征值：其特征多项式为
$(-X)^n$，由凯莱–哈密顿定理（\cref{ch:b2:reduction}）得 $u^n
= 0$，即幂零。
\end{solution}

\begin{solution}{exo:b2:linalg:11}
记该行列式为 $V(a_0, \dots, a_n)$，对 $n$ 作归纳；
以 $V(a_0) = 1$ 为起点。固定 $a_0, \dots, a_{n-1}$，视 $D(T) =
V(a_0, \dots, a_{n-1}, T)$，即末列为 $(1,
T, \dots, T^n)$：按该列展开，$D$ 是多项式，
次数 $\leq n$（关于 $T$），其 $T^n$ 项的系数为子式
$V(a_0, \dots, a_{n-1})$。先设 $a_0, \dots,
a_{n-1}$ 两两互异。对每个 $T = a_i$（$i < n$），有两列相同，
故 $D(a_i) = 0$。由 $n$ 个互异根，且次数
$\leq n$，于是
\[
D(T) = V(a_0, \dots, a_{n-1}) \prod_{i=0}^{n-1}(T - a_i),
\]
令 $T = a_n$，结合归纳假设，即得乘积
公式。若 $a_0, \dots, a_{n-1}$ 中有两个相等，则两边都
为 $0$（有重复列，亦有重复因子），
公式平凡成立。
\end{solution}

\begin{solution}{exo:b2:linalg:12}
\emph{$D$ 可逆。}右乘分块矩阵
$T = \begin{psmallmatrix} I & 0\\ -D^{-1}C & I\end{psmallmatrix}$,
该矩阵为对角块为单位阵的分块三角阵，故 $\det T = 1$（其
行列式由置换公式可知，只取对角线上的
块 --- 即\cref{exo:b2:linalg:2}中的分块规则）：
\[
\begin{pmatrix} A & B\\ C & D\end{pmatrix} T
= \begin{pmatrix} A - BD^{-1}C & B\\ C - DD^{-1}C & D\end{pmatrix}
= \begin{pmatrix} A - BD^{-1}C & B\\ 0 & D\end{pmatrix},
\]
其行列式为 $\det(A - BD^{-1}C)\det D = \det\bigl((A -
BD^{-1}C)D\bigr) = \det(AD - BD^{-1}CD)$。由于 $CD = DC$，
$BD^{-1}CD = BC$：行列式为 $\det(AD - BC)$。

\emph{一般情形 $D$。}令 $D_t = D + tI$，则仍有 $CD_t = D_tC$。
二者
\[
f(t) = \det\begin{pmatrix} A & B\\ C & D_t\end{pmatrix}
\qquad\text{与}\qquad
g(t) = \det(AD_t - BC)
\]
均为 $t$ 的多项式函数。多项式 $\det(D + tI)$ 是
首一的，次数为 $n$，故至多有 $n$ 个根：除有限多个
$t$ 之外，$D_t$ 可逆，且 $f(t) = g(t)$
由前一情形可得。无限域上的两个多项式在
无穷多点处取值相同，则二者相等：$f = g$，令 $t = 0$ 即得结论。
\end{solution}

\begin{solution}{pb:b2:linalg:1}
\textbf{1.} $\Phi$ 是线性映射，且 $\ker\Phi = \bigcap_i
\ker\varphi_i$（$p$ 元组为零，当且仅当每个分量为零）。对
坐标线性型 $\varepsilon_i$（$K^p$ 上的）：$\Phi^{\mathsf
T}(\varepsilon_i) = \varepsilon_i \circ \Phi = \varphi_i$，于是
$\operatorname{im}\Phi^{\mathsf T} \supseteq
\operatorname{Vect}(\varphi_i)$；反之
$\operatorname{im}\Phi^{\mathsf T}$ 由
$\Phi^{\mathsf T}(\varepsilon_i)$ 张成（$\varepsilon_i$ 张成
$(K^p)^*$）。故 $\operatorname{rk}\Phi = \operatorname{rk}
\Phi^{\mathsf T} = \dim\operatorname{Vect}(\varphi_1, \dots,
\varphi_p) =: r$（见\cref{prop:b2:linalg:transposerank}），
由秩--零化度定理可得 $\dim\bigcap_i\ker\varphi_i = n - r$。

\textbf{2.} （$\Leftarrow$）取一极大线性无关子族，不妨记为
$\varphi_1, \dots, \varphi_r$，它们张成同一空间（故
假设仍为 $\bigcap_{i \leq r}\ker\varphi_i \subseteq
\ker\varphi$：对全体 $i$ 的交与
$i \leq r$ 的交相等，因被舍去的每个线性型都是余下者的线性组合）。映射
$\Psi = (\varphi_1, \dots, \varphi_r) \colon E \to K^r$ 为
满射（由第 1 问，其秩为 $r$）。若 $\Psi(x) = \Psi(y)$
则 $x - y \in \ker\Psi \subseteq \ker\varphi$，故 $\varphi(x) =
\varphi(y)$：$\varphi$ 可分解为 $\varphi = \lambda \circ \Psi$
其中 $\lambda \colon K^r \to K$ 良定义；$\lambda$ 是线性的，
因 $\Psi$ 线性且满射（对 $t = \Psi(x)$，$t' =
\Psi(x')$: $\lambda(t + \alpha t') = \varphi(x + \alpha x') =
\lambda(t) + \alpha\lambda(t')$）。记 $\lambda = \sum c_i
\varepsilon_i$: $\varphi = \sum_{i \leq r} c_i\varphi_i$.
（$\Rightarrow$）若 $\varphi = \sum c_i \varphi_i$，则任一 $x$
使每个 $\varphi_i$ 为零，亦必使 $\varphi$ 为零。

\textbf{3.} 由第 1 问，$\dim\bigcap\ker\varphi_i = n - r$
其中 $r = \dim\operatorname{Vect}(\varphi_i) \leq p$，且 $r = p$
当且仅当该族线性无关。设 $F$ 为余维数为 $p$ 的子空间，其
零化子的维数为 $p$
（见\cref{thm:b2:linalg:annihilator}）；若 $(\varphi_1, \dots,
\varphi_p)$ 为 $F^\circ$ 的一组基，则 $F = \bigcap_i\ker\varphi_i$（
即恢复公式）。若线性型更少，则 $q$ 个超平面的交
维数 $\geq n - q > n - p$（由第 1 问）。

\textbf{4.} 计算 $\ker\varphi_1 \cap \ker\varphi_2$：由 $x
+ y - z = 0$ 与 $y + z - t = 0$，以 $(y, z)$ 参数化：$x = z
- y$，$t = y + z$，得向量构成的平面 $(z - y,\; y,\;
z,\; y + z)$。其上，$\psi = x + 2y - t = (z - y) + 2y - (y + z)
= 0$：由分解引理得 $\psi \in
\operatorname{Vect}(\varphi_1, \varphi_2)$ --- 事实上 $\psi =
\varphi_1 + \varphi_2$。但 $\psi' = x + y + t = (z - y) + y +
(y + z) = y + 2z$ 在其上并非恒为零（$y = 1, z = 0$
给出 $1$）：$\psi' \notin \operatorname{Vect}(\varphi_1,
\varphi_2)$.

\textbf{5.} $3$ 维空间上的三个线性型：只需
验证线性无关即可。若 $a\psi_0 + b\psi_1 + c\psi_2 = 0$，代入 $1, X,
X^2$: $a + b + c = 0$, $b + \frac c2 = 0$, $b + \frac c3 = 0$;
后两式相减得 $c = 0$，进而 $b = 0$，$a = 0$。
前对偶基：记 $P = \alpha + \beta X + \gamma X^2$，并解
方程 $\psi_i(P_j) = \delta_{ij}$（$P(0) = \alpha$，$P(1) =
\alpha + \beta + \gamma$, $\int_0^1 P = \alpha + \frac\beta2 +
\frac\gamma3$):
\[
P_0 = 1 - 4X + 3X^2, \qquad
P_1 = -2X + 3X^2, \qquad
P_2 = 6X - 6X^2 .
\]
（验证，例如：$\int_0^1 P_2 = 3 - 2 = 1$，$P_2(0) = P_2(1) = 0$。）
插值问题可由
前对偶基下的坐标求解：
\[
P = 1\cdot P_0 + 2\cdot P_1 + \tfrac32\, P_2 = 1 + X
\]
（$X$ 的系数 $-4 - 4 + 9 = 1$，$X^2$ 的系数 $3 + 6 - 9 =
0$）；事实上 $P(0) = 1$，$P(1) = 2$，$\int_0^1 P = \frac32$。

\textbf{6.} 线性性：对每个 $\varphi$，$J(x + \alpha
y)(\varphi) = \varphi(x + \alpha y) = J(x)(\varphi) + \alpha
J(y)(\varphi)$，即 $J(x + \alpha y) = J(x) + \alpha J(y)$。单射性：若 $x \neq 0$，
将 $x = e_1$ 扩充为一组基；坐标线性型 $e_1^*$ 有
$J(x)(e_1^*) = 1 \neq 0$。由于 $\dim E^{**} = \dim E^* = \dim
E$，故单射即双射。

\textbf{7.} 包含关系：对 $x \in F$ 与 $\varphi \in F^\circ$，
$J(x)(\varphi) = \varphi(x) = 0$，故 $J(F) \subseteq
F^{\circ\circ}$。比较维数
（两次应用\cref{thm:b2:linalg:annihilator}）：
\[
\dim F^{\circ\circ} = \dim E^* - \dim F^\circ
= n - (n - \dim F) = \dim F = \dim J(F),
\]
因 $J$ 为单射，故 $J(F) = F^{\circ\circ}$。

\textbf{8.} 第一个恒等式：$\varphi$ 零化 $F + G$，当且仅当它
同时零化 $F$ 与 $G$（零化和，当且仅当零化各个加项）：
$(F+G)^\circ = F^\circ \cap G^\circ$。第二个为包含关系
$F^\circ + G^\circ \subseteq (F \cap G)^\circ$ 显然成立（每个
加项都零化 $F \cap G$）。计算维数时用到第一个恒等式
与格拉斯曼公式：
\[
\dim(F^\circ + G^\circ) = \dim F^\circ + \dim G^\circ -
\dim(F^\circ \cap G^\circ)
= (n - \dim F) + (n - \dim G) - \bigl(n - \dim(F +
G)\bigr),
\]
由 $E$ 中的格拉斯曼公式，上式等于 $n - \dim(F \cap G) = \dim(F
\cap G)^\circ$，等式成立。

\textbf{9.} 由引理，$\ker\psi \subseteq \ker\varphi$ 且
$p = 1$ 时得 $\varphi \in \operatorname{Vect}(\psi)$，且
$\varphi \neq 0$ 使标量非零。直接地，取 $x_0$
使 $\psi(x_0) \neq 0$；每个 $x$ 可写为 $x = \bigl(x -
\frac{\psi(x)}{\psi(x_0)}x_0\bigr) + \frac{\psi(x)}{\psi(x_0)}
x_0$，其中第一项属于 $\ker\psi = \ker\varphi$；
对其应用 $\varphi$，得 $\varphi(x) = \frac{\varphi(x_0)}{\psi(x_0)}\psi(x)$。

\textbf{10.} 取 $(\varphi_1^*, \dots,
\varphi_n^*)$ 为 $(\varphi_1, \dots, \varphi_n)$ 在 $E^{**}$ 中的对偶基，
（\cref{def:b2:linalg:dual}用于 $E^*$），并令 $u_j =
J^{-1}(\varphi_j^*)$：它构成 $E$ 的基（$J$ 为同构，
见第 6 题），且 $\varphi_i(u_j) = J(u_j)(\varphi_i) =
\varphi_j^*(\varphi_i) = \delta_{ij}$。唯一性：条件
$\varphi_i(u_j) = \delta_{ij}$ 确定 $J(u_j)$ 在
基 $(\varphi_i)$ 上的值，从而确定 $u_j$。

\textbf{11.} 线性性：$(u + \alpha v)^{\mathsf T}\psi = \psi
\circ (u + \alpha v) = u^{\mathsf T}\psi + \alpha\, v^{\mathsf
T}\psi$。单射性：若 $u \neq 0$，取 $x$ 使 $u(x) \neq 0$，
再取 $\psi$ 使 $\psi(u(x)) \neq 0$（利用第 6 题的坐标形式
技巧）：$u^{\mathsf T}\psi \neq 0$。空间 $\mathcal{L}(E,F)$
与 $\mathcal{L}(F^*, E^*)$ 的维数均为 $\dim E \dim F$，
故为双射。若 $u$ 可逆，则由反序律 $(vu)^{\mathsf
T} = u^{\mathsf T}v^{\mathsf T}$，于是 $u^{\mathsf
T}(u^{-1})^{\mathsf T} = (u^{-1}u)^{\mathsf T} =
\mathrm{id}_{E^*}$，且 $(u^{-1})^{\mathsf T}u^{\mathsf T} =
(uu^{-1})^{\mathsf T} = \mathrm{id}_{F^*}$，故 $(u^{\mathsf
T})^{-1} = (u^{-1})^{\mathsf T}$。

\textbf{12.} 对 $x \in E$ 与 $\psi \in F^*$：
\[
\bigl(u^{\mathsf T\mathsf T}(J_E x)\bigr)(\psi)
= (J_E x)\bigl(u^{\mathsf T}\psi\bigr)
= (u^{\mathsf T}\psi)(x)
= \psi\bigl(u(x)\bigr)
= \bigl(J_F(u(x))\bigr)(\psi).
\]
由 $\psi$ 的任意性，$u^{\mathsf T\mathsf T} \circ J_E = J_F
\circ u$。

\textbf{13.} 由\cref{prop:b2:linalg:transposerank}可知：$\ker
u^{\mathsf T} = (\operatorname{im} u)^\circ$，故 $u$ 满射
$\iff \operatorname{im} u = F \iff (\operatorname{im}u)^\circ =
\{0\}$（\cref{thm:b2:linalg:annihilator}）$\iff u^{\mathsf T}$
为单射。又 $\operatorname{im} u^{\mathsf T} = (\ker
u)^\circ$，故 $u$ 单射 $\iff \ker u = \{0\} \iff (\ker
u)^\circ = E^*$ $\iff u^{\mathsf T}$ 满射。

\textbf{14.} 若 $u(F) \subseteq F$ 且 $\varphi \in F^\circ$：
$(u^{\mathsf T}\varphi)(x) = \varphi(u(x)) = 0$，对 $x \in F$，故
$u^{\mathsf T}\varphi \in F^\circ$。反之，若 $u(F)
\not\subseteq F$，取 $x \in F$ 使 $u(x) \notin F$；由
\cref{thm:b2:linalg:annihilator}的恢复公式，存在
$\varphi \in F^\circ$ 使 $\varphi(u(x)) \neq 0$。于是
$(u^{\mathsf T}\varphi)(x) \neq 0$，尽管 $x \in F$，故
$u^{\mathsf T}\varphi \notin F^\circ$，即 $F^\circ$ 不稳定。

\textbf{15.} $u^{\mathsf T} - \lambda\,\mathrm{id}_{E^*} = (u -
\lambda\,\mathrm{id}_E)^{\mathsf T}$（转置是线性的，且
$\mathrm{id}^{\mathsf T} = \mathrm{id}$），故其核为
$(\operatorname{im}(u - \lambda\,\mathrm{id}))^\circ$
（\cref{prop:b2:linalg:transposerank}），其维数为
\[
n - \operatorname{rk}(u - \lambda\,\mathrm{id})
= \dim\ker(u - \lambda\,\mathrm{id})
\]
由秩--零化度定理。特别地，一个核非零当且仅当
另一个核非零：特征值相同，几何重数也相同。

\textbf{16.} 包含关系：若 $b = u(x)$ 且 $u^{\mathsf T}\psi =
0$，则 $\psi(b) = \psi(u(x)) = (u^{\mathsf T}\psi)(x) = 0$，故
$\operatorname{im} u \subseteq (\ker u^{\mathsf T})_\circ$。
比较维数：对子空间 $S \subseteq F^*$，$S_\circ =
J_F^{-1}(S^\circ)$（展开即得：$y \in S_\circ$ 当且仅当每个 $\psi \in
S$ 零化 $y$，亦即 $J_F(y) \in S^\circ$），故 $\dim S_\circ = \dim
F - \dim S$。取 $S = \ker u^{\mathsf T}$：
\[
\dim(\ker u^{\mathsf T})_\circ
= \dim F - \dim\ker u^{\mathsf T}
= \operatorname{rk} u^{\mathsf T} = \operatorname{rk} u :
\]
维数相等，故 $\operatorname{im} u = (\ker
u^{\mathsf T})_\circ$。换言之，$b \in \operatorname{im} u$ 当且仅当
$\psi(b) = 0$ 对每个 $\psi$ 成立，只要 $u^{\mathsf T}\psi = 0$ ---
即弗雷德霍姆择一性。

\textbf{17.} 将 $(K^m)^*$ 与 $K^m$ 通过 $y \mapsto \psi_y$ 等同，
$\psi_y(v) = y^{\mathsf T}v$，于是 $(u^{\mathsf T}\psi_y)(x) =
y^{\mathsf T}Ax = (A^{\mathsf T}y)^{\mathsf T}x$，故 $u^{\mathsf
T}\psi_y = \psi_{A^{\mathsf T}y}$：转置映射即为
转置矩阵。\emph{至多一个解：}若 $Ax = b$ 且
$A^{\mathsf T}y = 0$，则 $y^{\mathsf T}b = y^{\mathsf T}Ax =
(A^{\mathsf T}y)^{\mathsf T}x = 0 \neq 1$。\emph{至少一个解：}若
(i) 不成立，则由第 16 题给出 $\psi_y$，满足 $A^{\mathsf T}y
= 0$ 且 $y^{\mathsf T}b \neq 0$；对 $y$ 作数乘，使之为 $1$。

\textbf{18.} $A = \begin{psmallmatrix} 1 & 1 & 0\\ 0 & 1 & 1\\
1 & 2 & 1\end{psmallmatrix}$（第三行 = 第一行 + 第二行，故 $A$
是奇异的）。求解 $A^{\mathsf T}y = 0$：$y_1 + y_3 = 0$，$y_1
+ y_2 + 2y_3 = 0$，$y_2 + y_3 = 0$ 得 $y_1 = y_2 = -y_3$：
由 $y = (1, 1, -1)$ 张成的直线。弗雷德霍姆：可解当且仅当
$y^{\mathsf T}b = b_1 + b_2 - b_3 = 0$，即 $b_3 = b_1 + b_2$
--- 显然这正是正确的条件，因为第三个方程是
前两个方程之和。

\textbf{19.} $L$ 的矩阵对角线上为 $1$，
而 $-\frac12$ 位于 $(k, k\pm1)$（模 $n$）处，该矩阵对称，故
$L^{\mathsf T} = L$（依第 17 题的等同）。
\emph{核：}若 $Lx = 0$，则每个 $x_k = \frac12(x_{k-1} +
x_{k+1})$。设 $k_0$ 使 $x_k$ 最大；两个邻点的
算术平均值（二者均 $\leq x_{k_0}$）等于 $x_{k_0}$ 仅当
两者都等于 $x_{k_0}$；沿圈传递，$x$ 为常数。
反之，常值向量显然属于核。因此 $\ker L^{\mathsf T} = \ker L
= \R(1, \dots, 1)$，弗雷德霍姆择一性表述为：$Lx = b$
可解当且仅当 $(1,\dots,1)^{\mathsf T} b = \sum_k b_k = 0$ ---
即离散相容性条件：环上的“热分布”
可由势函数实现，当且仅当其总通量
为零。

\textbf{20.} 若 $A$ 反对称，$S$ 对称：
\[
\operatorname{tr}(AS) = \operatorname{tr}\bigl((AS)^{\mathsf
T}\bigr) = \operatorname{tr}(S^{\mathsf T}A^{\mathsf T}) =
-\operatorname{tr}(SA) = -\operatorname{tr}(AS),
\]
故 $2\operatorname{tr}(AS) = 0$，且（$\operatorname{char} K \neq
2$）时 $\operatorname{tr}(AS) = 0$：$\mathcal{A}_n \subseteq
\mathcal{S}_n^\circ$（将对偶与矩阵等同）。
维数：$\dim\mathcal{S}_n^\circ = n^2 - \frac{n(n+1)}2 =
\frac{n(n-1)}2 = \dim\mathcal{A}_n$：等号成立。互换角色
（计算相同），$\mathcal{A}_n^\circ = \mathcal{S}_n$。

\textbf{21.} $\operatorname{tr}(I_nM) = \operatorname{tr} M = 0$
对 $M \in \mathfrak{sl}_n$，直线 $KI_n$ 落在
零化子中，其维数为 $n^2 - (n^2 - 1) = 1$：等号成立。
经同构 $A \mapsto
\operatorname{tr}(A\,\cdot)$ 转译：在
$\mathfrak{sl}_n$ 上为零的形式是 $\operatorname{tr}(\lambda I_n\,\cdot) =
\lambda\operatorname{tr}$。

\textbf{22.} 设 $M \in \mathcal{M}_n(K)$。多项式 $t
\mapsto \det(M - tI)$ 非零且次数为 $n$，故至多有
$n$ 个根；$K$ 特征为 $0$，故为无限域，取
一个不是根的 $\lambda \neq 0$。于是 $M = (M - \lambda I) +
\lambda I$ 将 $M$ 写成两个可逆矩阵之和。

\textbf{23.} \emph{第一步：}对可逆 $P$ 及任意
$X$，对 $M = XP$ 应用不变性：$t(P(XP)P^{-1}) = t(XP)$，
即 $t(PX) = t(XP)$。\emph{第二步：}固定 $X$，两边
$t(BX) = t(XB)$ 关于 $B$ 均为线性，且在可逆 $B$ 上相等；
由第 22 题，每个 $B$ 都是两个可逆矩阵之和，故二者
处处相等。\emph{第三步：}$t$ 将每个换位子 $XB
- BX$ 映为零；这些换位子张成 $\mathfrak{sl}_n$（在
\cref{prop:b2:linalg:trace}的证明中已给出），故 $t$ 在
$\mathfrak{sl}_n$ 上为零，由第 21 题得 $t =
c\operatorname{tr}$。（反之，每个 $c\operatorname{tr}$ 都是相似不变的：迹
恰是\emph{唯一}的线性相似
不变量。）

\textbf{24.} 若 $\operatorname{rk} u = r' \leq r$：取
$(f_1, \dots, f_{r'})$ 为 $\operatorname{im} u$ 的一组基，并写 $u(x)
= \sum_{i=1}^{r'} \psi_i(x) f_i$；$\psi_i(x)$ 是
$u(x)$ 的坐标，关于 $x$ 线性（$u$ 与坐标
形式的复合），故 $u$ 为 $r' \leq r$ 个秩 $\leq1$ 的映射之和（
补零）。反之，若 $u = \sum_{i=1}^r \psi_i(\cdot)f_i$，
则 $\operatorname{im} u \subseteq \operatorname{Vect}(f_1,
\dots, f_r)$，故 $\operatorname{rk} u \leq r$。次可加性：写
$u$ 为含 $\operatorname{rk} u$ 项之和，$v$ 为含
$\operatorname{rk} v$ 项；故和式共含 $\operatorname{rk} u +
\operatorname{rk} v$ 项，所以 $\operatorname{rk}(u + v) \leq
\operatorname{rk} u + \operatorname{rk} v$。

\textbf{25.} 对应关系：子空间 $F$ 对应于
$F^\circ$，二者维数互补
（见\cref{thm:b2:linalg:annihilator}），再由双对偶返回
（第 6--7 题）；和与交互换（第 8 题）；
映射 $u$ 对应 $u^{\mathsf T}$，且 $\ker u^{\mathsf T}
= (\operatorname{im}u)^\circ$，$\operatorname{im}u^{\mathsf T} =
(\ker u)^\circ$，秩相等，单射性与满射性
互换，对应的稳定子空间与
特征值（第 11--15 题）；方程 $u(x) = b$
可解当且仅当 $b$ 正交于 $\ker u^{\mathsf T}$
（第 16--19 题）；在 $\mathcal{M}_n$ 上迹配对
具体实现整个对应关系，以迹为
唯一的线性相似不变量（第 20--23 题），而秩
则为分解为初等张量的最小长度
（第 24 题）。在无限维情形下，维数
计数失效，代之以像上的\emph{闭性}假设，以及
完备性 --- 在希尔伯特空间上，这便成为
里斯表示定理，以及紧算子的
弗雷德霍姆理论；这些内容将在本科三年级分册中严格证明。
\end{solution}
